\section{Methodology}
\label{sec:method}

\paragraph{Controlled tokenizers (Track A).} We train byte-level BPE tokenizers
\citep{radford2019gpt2} with the Hugging Face \texttt{tokenizers} implementation
\citep{huggingface2026tokenizers}, holding fixed: a byte-level pre-tokenizer and
matching decoder, a full byte-level initial alphabet (so ordinary text never
produces an unknown token), Unicode NFC normalization without lowercasing or
accent stripping, the special-token set, and a minimum merge frequency of two.
The only variable is the training corpus (General, Azure, Mixed). Vocabulary sizes
are $8{,}192$, $16{,}384$, and $32{,}768$ (special tokens count toward the size),
giving nine controlled tokenizers. Sampling, splitting, and the bootstrap are
driven by a single seed ($20{,}260{,}709$). Corpora are sampled deterministically
to exactly $33{,}554{,}432$ bytes each; the Mixed corpus interleaves Azure and
General documents 50/50 by bytes. Table~\ref{tab:matrix} lists the design.

\begin{table}[t]
\centering
\caption{Controlled experimental design. Only the training corpus is manipulated;
all other factors are held fixed within a comparison.}
\label{tab:matrix}
\small

\begin{tabular}{lll}
\toprule
\textbf{Factor} & \textbf{Levels} & \textbf{Held constant} \\
\midrule
Training corpus (manipulated) & General, Azure, Mixed 50/50 & --- \\
Vocabulary size & 8{,}192 / 16{,}384 / 32{,}768 & within a comparison \\
Algorithm & byte-level BPE & implementation, pre-tokenizer, decoder \\
Training budget & 33,554,432 UTF-8 bytes & byte-matched across corpora \\
Normalization & Unicode NFC & no lowercasing, no accent stripping \\
Seed & 20260709& sampling, splitting, bootstrap \\
\bottomrule
\end{tabular}

\end{table}

\paragraph{Observational baselines (Track B).} We also encode the evaluation sets
with production \texttt{tiktoken} encodings (\texttt{r50k\_base},
\texttt{cl100k\_base}, \texttt{o200k\_base}) \citep{openai2026tiktoken}. These
differ from Track A and from one another in training data, vocabulary size,
normalization, and intended model, so they are reported separately and are not a
controlled comparison.

\paragraph{Primary metric.} For a tokenizer $T$ and evaluation corpus $D$,
\begin{equation}
\mathrm{tpk}(T, D) = 1000 \cdot \frac{\mathrm{tokens}(T, D)}{\mathrm{bytes}_{\text{UTF-8}}(D)},
\end{equation}
tokens per $1{,}000$ UTF-8 bytes (lower is more compression), reported both
\emph{micro} (over all bytes) and \emph{macro} (mean of per-document ratios).
The primary effect size is the relative reduction of a specialized tokenizer $A$
versus the matched General tokenizer $G$,
\begin{equation}
\Delta(A \text{ vs } G, D) = 100 \cdot \frac{\mathrm{tpk}(G, D) - \mathrm{tpk}(A, D)}{\mathrm{tpk}(G, D)},
\end{equation}
where \emph{positive} $\Delta$ means the specialized tokenizer uses \emph{fewer}
tokens. Secondary metrics include characters/bytes per token, tokens per
whitespace word, identifier fragmentation, vocabulary utilization, exact
round-trip and unknown-token rates, and CPU throughput.

\paragraph{Statistics.} Token counts are deterministic; uncertainty comes from the
sampled documents. We use a paired bootstrap over documents ($10{,}000$ resamples;
the same documents resampled for both tokenizers), reporting 95\% confidence
intervals for $\Delta$, a document-level win rate, and stratification by
evaluation set, content type, and repository.

\paragraph{Runtime.} CPU tokenizer throughput is measured separately from
compression, with a fixed thread count, warm-up, five measured rounds, randomized
tokenizer order per round, identical documents and batches, and serialized
execution, on an Apple M3 Pro (12 cores, macOS). This is CPU tokenizer throughput,
not model-serving throughput.

\paragraph{Verification.} An automated gate re-checks byte-matching, leakage
controls, exact round-trip and unknown-token counts, totals reconciliation,
provenance completeness, artifact-distribution policy, and source pinning without
retraining. Every number in this paper is additionally reconciled against its
source file by a claim registry (Appendix~\ref{app:verify}).
